\documentclass[../../../include/open-logic-section]{subfiles}

\begin{document}

\olfileid{sth}{story}{blv}

\olsection{ملحق: قانون فرغه الأساسي الخامس}

شرحنا في~\olref[rus]{sec} أن راسل صاغ مفارقته بوصفها مشكلة للنسق
الذي عرضه فرغه في~\emph{Grundgesetze}. ولم يتضمن نسق فرغه صياغة
مباشرة للاستيعاب الساذج. ولذلك سنشرح بإيجاز شديد في هذا الملحق ما
\emph{تضمنه} نسق فرغه، وكيف يتصل بالاستيعاب الساذج وبمفارقة راسل.

نسق فرغه \emph{من الرتبة الثانية}، وقد صُمم لصوغ مفهوم
\emph{امتداد مفهوم}.\footnote{يحاول فرغه، على وجه الدقة، إضفاء الصبغة
  الصورية على مفهوم أعم هو «مدى قيم» دالة. وامتدادات المفاهيم حالة خاصة
  من ذلك المفهوم الأعم. للتفاصيل، انظر
  \citet[pp.\ 8--9]{Heck2012}.} وبترميز مستلهم من فرغه، سنكتب
$\fregeext{x}{F(x)}$ للدلالة على \emph{امتداد المفهوم~$F$}. وهذه
أداة تأخذ \emph{محمولًا}، «$F$»، وتحوله إلى \emph{حد} (من الرتبة
الأولى)، «$\fregeext{x}{F(x)}$». وباستعمال هذه الأداة، قدم فرغه
\emph{التعريف} الآتي للانتماء:
\[
	a \in b \liff_\text{df} \exists G(b = \fregeext{x}{G(x)} \land Ga)
\]
أي، على وجه التقريب، إن~$a \in b$ إذا وفقط إذا اندرج~$a$ تحت مفهوم
امتداده~$b$. (لاحظ أن المكمّم «$\exists G$» من الرتبة الثانية.)
وقد تبنى فرغه أيضًا المبدأ الآتي، المعروف بـ\emph{القانون الأساسي
الخامس}:
$$\fregeext{x}{F(x)} = \fregeext{x}{G(x)} \liff \forall x (Fx \liff Gx)$$
أي، على وجه التقريب، إن للمفهومين امتدادًا واحدًا إذا وفقط إذا كانا
متساويي الامتداد. (ومرة أخرى، يقع كل من «$F$» و«$G$» في موضع المحمول.)
ويربط الآن مبدأ بسيط الانتماء بتحقق الخاصية:

\begin{lem}[في \emph{Grundgesetze}]\ollabel{lem:Fregeextensions}
$\forall F \forall a(a \in \fregeext{x}{F(x)} \liff Fa)$
\end{lem}

\begin{proof}
ثبّت~$F$ و$a$. يكون~$a \in \fregeext{x}{F(x)}$ إذا وفقط إذا كان
$\exists G(\fregeext{x}{F(x)} = \fregeext{x}{G(x)} \land Ga)$
(بتعريف الانتماء)، وإذا وفقط إذا كان
$\exists G(\forall x(Fx \liff Gx) \land Ga)$
(بالقانون الأساسي الخامس)، وإذا وفقط إذا كان~$Fa$
(بالمنطق الأولي من الرتبة الثانية).
\end{proof}

وينتج من ذلك الاستيعاب الساذج على الفور تقريبًا:

\begin{lem}[في \emph{Grundgesetze}.]
$\forall F \exists s \forall a (a \in s \liff Fa)$
\end{lem}

\begin{proof}
ثبّت~$F$؛ ينتج الآن من~\olref{lem:Fregeextensions} أن
$\forall a (a \in \fregeext{x}{F(x)} \liff Fa)$؛ ومن ثم
$\exists s\forall a(a \in s \liff Fa)$ بالتعميم الوجودي. وتنتج النتيجة
لأن~$F$ كانت اعتباطية.
\end{proof}

وتنتج مفارقة راسل بأخذ~$F$ كما تعطيها
$\forall x(Fx \liff x \notin x)$.

\end{document}
